 \chapter{数到 20}\label{ch:g1:counting}

 有多少颗弹珠？数数回答的是
数学中最基本的问题。本章我们数一数较少的物品，
写出从 $0$ 到 $20$ 的数，并进行比较：谁的更多？

 \section{数物体}

 \begin{method}[怎样数数]\label{met:g1:counting:how}
\begin{enumerate}
  \item 每个物体只摸一次，并按顺序说出数字：一，
        二、三、\dots；
  \item 不要漏掉任何物体，也不要重复摸同一个（把数过的物体移开，
        或者把它们划掉）；
  \item \emph{最后}说出的数字就是物体的个数。
\end{enumerate}
数的顺序不影响结果：从左往右数，
还是从右往左数，$7$ 颗弹珠都是 $7$ 颗。
\end{method}

\begin{omfigure}
\begin{tikzpicture}[scale=0.75]
  \foreach \i/\n in {0/1,1/2,2/3,3/4,4/5} {
    \fill[omDef] (\i*1.1,0) circle (0.33);
    \node[above, font=\small] at (\i*1.1,0.45) {\n};
  }
  \node[right, font=\small] at (5.2,0) {五颗弹珠：共有 $5$ 颗};
\end{tikzpicture}

{\small 数数：每颗弹珠对应一个数，最后一个数字
就是答案。}
\end{omfigure}

\begin{definition}[零]\label{def:g1:counting:zero}
当没有东西可数 --- 比如一个空盒子 --- 时，数是
\emph{零}\index{零}，记作 $0$。
\end{definition}

 \section{20 以内的数}

\begin{example}[读数与写数]\label{ex:g1:counting:writing}
按顺序写出从 $0$ 到 $20$ 的数：
\[
0,\ 1,\ 2,\ 3,\ 4,\ 5,\ 6,\ 7,\ 8,\ 9,\ 10,
\]
\[
11,\ 12,\ 13,\ 14,\ 15,\ 16,\ 17,\ 18,\ 19,\ 20 .
\]
$9$ 之后，数要用\emph{两}位数来写：$13$ 是一个十和
$3$ 个一；$20$ 是两个十。（
十位将在\cref{ch:g1:tens}中学习。）
\end{example}

\begin{omfigure}
\begin{tikzpicture}[scale=0.62]
  \draw[-Stealth] (-0.4,0) -- (20.8,0);
  \foreach \i in {0,...,20} {
    \draw (\i,-0.12) -- (\i,0.12);
    \node[below, font=\footnotesize] at (\i,-0.15) {\i};
  }
  \fill[omThm] (7,0) circle (2.6pt);
  \fill[omMeth] (13,0) circle (2.6pt);
\end{tikzpicture}

{\small 从 $0$ 到 $20$ 的数轴：每个数都有自己的位置。
向右走，数越来越大。圆点标出了 $7$ 和 $13$。}
\end{omfigure}

\begin{example}[前一个，后一个]\label{ex:g1:counting:neighbors}
在数轴上，每个数都有两个邻居：$14$ 前面是
$13$，后面是 $15$。倒着数 --- $10$、$9$、$8$、
\dots\ --- 就是向左走。
\end{example}

 \section{比较}

\begin{method}[谁的更多？]\label{met:g1:counting:compare}
比较两组物体时，可以一一配对：每组各取一个，
并排放在一起。
\begin{enumerate}
  \item 若某一组中有物体没有配到对，
        则这一组\emph{更多}；
  \item 若两边都没有剩下，
        则两组\emph{一样多}。
\end{enumerate}
对于数，不需要配对：在数轴上更靠右的那个数
更大 --- 配对和数轴的结论总是一致。
记作 $12 > 9$（$12$ 更多）或 $9 < 12$（$9$ 更少）。
\end{method}

\begin{omfigure}
\begin{tikzpicture}[scale=0.62]
  \foreach \i in {0,...,6} \fill[omDef] (\i*1.1,1.3) circle (0.3);
  \foreach \i in {0,...,4} \fill[omThm] (\i*1.1,0) circle (0.3);
  \foreach \i in {0,...,4}
    \draw[omExo, thick] (\i*1.1,0.3) -- (\i*1.1,1.0);
  \draw[omExo, thick] (5.5,1.3) circle (0.52);
  \draw[omExo, thick] (6.6,1.3) circle (0.52);
  \node[right, font=\small] at (7.5,0.65)
    {$7 > 5$：两颗蓝色弹珠没有配对};
\end{tikzpicture}

{\small 通过配对比较：$7$ 颗蓝色弹珠与 $5$ 颗红色弹珠相比，
--- 圈出的两颗蓝色弹珠剩了下来，所以 $7$ 比 $5$ 多。}
\end{omfigure}

\begin{example}[第一、第二、第三]\label{ex:g1:counting:ordinal}
赛跑时，选手按先后顺序到达：\emph{第一}、
\emph{第二}、\emph{第三}、\emph{第四}\,\dots\ 这些
词说明的是\emph{位置}，而不是数量：第三名选手
是一名选手，排在第 $3$ 位。
\end{example}

\section{练习}

\begin{exercise}[$\star$]\label{exo:g1:counting:1}
数一数你名字中有几个字母。
数一数两只手有几根手指，数一数两只猫有几条腿。
\end{exercise}

\begin{exercise}[$\star$]\label{exo:g1:counting:2}
用数字写出：七；\ 十二；\ 二十；\ 零。用汉字写出：
$4$；\ $15$；\ $18$。
\end{exercise}

\begin{exercise}[$\star$]\label{exo:g1:counting:3}
接着数下去：$8,\ 9,\ 10,\ ?,\ ?,\ ?$。倒着数：
$15,\ 14,\ 13,\ ?,\ ?,\ ?$。
\end{exercise}

\begin{exercise}[$\star$]\label{exo:g1:counting:4}
$10$ 的前一个数是哪一个？$16$ 的后一个数是哪一个？
$12$ 和 $14$ 之间的数是哪一个？
\end{exercise}

\begin{exercise}[$\star$]\label{exo:g1:counting:5}
照样子抄写，并用 $<$ 或 $>$ 填空：
\[
6 \;?\; 9, \qquad
14 \;?\; 8, \qquad
11 \;?\; 13, \qquad
20 \;?\; 19 .
\]
\end{exercise}

\begin{exercise}[$\star$]\label{exo:g1:counting:6}
画出 $6$ 个圆和 $8$ 个叉，通过配对，
哪一种更多？多几个？
\end{exercise}

\begin{exercise}[$\star$]\label{exo:g1:counting:7}
从小到大排列：$12$；\ $7$；\ $17$；\ $2$；\ $10$。
\end{exercise}

\begin{exercise}[$\star$]\label{exo:g1:counting:8}
五个小朋友排成一队：爱丽丝、鲍勃、克莱奥、丹、艾玛。
谁排第二？谁排第四？丹排在第几位？
\end{exercise}

\begin{exercise}[$\star\star$]\label{exo:g1:counting:9}
我是一个介于 $10$ 和 $20$ 之间的数。
按数线数，我在 $13$ 之后第三步的位置。
我是谁？而在 $13$ 的\emph{之前}第三步的数是谁？
\end{exercise}
